% Teil II: durchgerechneter Weg, Teile a--c.

\subsection*{Gegeben und gesucht}
$a(t)$ steigt fuer $0\le t<T$ linear von $0$ auf $a_0$, faellt fuer
$T\le t<2T$ linear auf $0$ und ist fuer $t\ge2T$ null.
Anfangsbedingungen $v(0)=0$, $x(0)=0$. Gesucht: (a) $a(t)$ in den drei
Intervallen, (b) $v(t)$ und $x(t)$ fuer $t\ge0$, (c) $v_{\mathrm{end}}$.

\subsection*{Ansatz}
Die Beschleunigung ist in jedem Intervall eine Gerade; einmaliges bzw.
zweimaliges Integrieren liefert $v(t)$ bzw. $x(t)$. Die Randwerte werden
uemer die Anschlussbedingungen verknuepft.

\subsection*{Rechnung}
\textbf{a) Beschleunigung.}
\begin{equation}
  a(t)=\begin{cases}
    \dfrac{a_0}{T}\,t & 0\le t<T,\\[1.5ex]
    a_0\left(2-\dfrac{t}{T}\right) & T\le t<2T,\\[1.5ex]
    0 & t\ge2T.
  \end{cases}
\end{equation}

\textbf{b) Geschwindigkeit und Ort.} Fuer das erste Intervall:
\begin{equation}
  v(t)=\int_0^t \frac{a_0}{T}t'\,\mathrm dt'=\frac{a_0\,t^2}{2T},
  \qquad x(t)=\int_0^t v\,\mathrm dt'=\frac{a_0\,t^3}{6T},
\end{equation}
also $v(T)=a_0T/2$ und $x(T)=a_0T^2/6$.

Zweites Intervall mit $v(T)$ als Startwert:
\begin{equation}
  v(t)=v(T)+\int_T^t a_0\left(2-\frac{t'}{T}\right)\mathrm dt'
  = a_0\left(2t-\frac{t^2}{2T}-T\right),
\end{equation}
\begin{equation}
  x(t)=x(T)+\int_T^t v(t')\,\mathrm dt'
  = a_0\left(t^2-\frac{t^3}{6T}-Tt+\frac{T^2}{3}\right).
\end{equation}
Am rechten Rand $t=2T$: $v(2T)=a_0T$, $x(2T)=a_0T^2$.

Drittes Intervall ($a=0\Rightarrow v$ konstant):
\begin{equation}
  v(t)=a_0T,\qquad
  x(t)=a_0T^2+a_0T\,(t-2T)=a_0T\,(t-T)\;.
\end{equation}

\textbf{c) Endgeschwindigkeit.}
\begin{equation}
  \boxed{\;v_{\mathrm{end}}=v(2T)=a_0T\;}.
\end{equation}
Die Skizze zeigt $v(t)$: steigende Parabel, dann fallende Parabel, ab $2T$
eine Horizontale auf $a_0T$.

\begin{ergebnis}
$a(t)$ wie oben; $v(t)$ und $x(t)$ stueckweise mit den Anschluessen
$v(T)=\tfrac12a_0T$, $v(2T)=a_0T$; $v_{\mathrm{end}}=a_0T$.
\end{ergebnis}

\begin{kontrolle}
\textbf{Einheiten:} $a_0$ in \unit{\metre\per\second\squared}, $T$ in
\unit{\second} --- $a_0T$ damit in \unit{\metre\per\second}. \checkmark

\textbf{Querprobe:} Fuer eine Dreiecksspitze der Breite $2T$ und Hoehe
$a_0$ ist die Flaeche $\int a\,\mathrm dt=\frac12(2T)a_0=a_0T$, genau
$v_{\mathrm{end}}$. \checkmark

\textbf{Stetigkeit:} $v(T)=\tfrac12a_0T$ aus beiden angrenzenden
Formeln; keine Spruenge. \checkmark
\end{kontrolle}

\begin{merksatz}
Die Geschwindigkeit ist die Flaeche unter $a(t)$. Bei stueckweise
linearem $a(t)$ integriert man Intervall fuer Intervall und schliesst die
Randwerte an; die Enddgeschwindigkeit ist die gesamte Flaeche unter der
Beschleunigungskurve.
\end{merksatz}
